\chapter{Prompts utilizados}
\label{anexo:prompts}

\section{Criterio de documentación}

Este anexo documenta una selección representativa de los prompts utilizados durante el desarrollo. No se reproducen de forma literal: se presentan en una versión revisada, más clara y reutilizable, manteniendo el objetivo real de cada consulta. La decisión responde a dos motivos. Primero, muchos prompts originales formaban parte de conversaciones largas, con adjuntos, correcciones sucesivas y mensajes de seguimiento que no resultan legibles como evidencia académica. Segundo, el interés del anexo no es conservar cada frase escrita, sino explicar cómo se utilizó la IA generativa dentro del proceso de ingeniería.

Las respuestas de los asistentes se trataron como apoyo para generar hipótesis, ordenar alternativas y proponer comprobaciones. Ninguna decisión se incorporó únicamente por aparecer en una respuesta: se contrastó con Blueprints, CSV, informes de análisis, pruebas en Unreal Engine o revisión manual. Las consultas menores sobre uso básico del editor, avisos de LaTeX o formato no se incluyen salvo que ayudasen a reconstruir una decisión técnica.

\section{Material revisado}

La carpeta \path{INFO/CONVERSACIONES_IA} contiene 48 conversaciones exportadas en Markdown. La tabla~\ref{tab:resumen-prompts-ia} resume el material revisado.

\begin{table}[htbp]
\centering
\caption{Conversaciones de IA revisadas para el anexo de prompts.}
\label{tab:resumen-prompts-ia}
\footnotesize
\setlength{\tabcolsep}{4pt}
\renewcommand{\arraystretch}{0.92}
\begin{tabular}{L{2.2cm}rrL{7.2cm}}
\toprule
\textbf{Asistente} & \textbf{Ficheros} & \textbf{Bloques} & \textbf{Uso principal}\\
\midrule
ChatGPT & 13 & 209 & Depuración de Blueprints, scripts de análisis, modelos y memoria.\\
Claude & 28 & 292 & Análisis de entrenamientos, regresiones, stops, rotondas y selección de modelos.\\
Gemini & 7 & 268 & Apoyo inicial con Unreal Engine, organización y consultas puntuales.\\
\bottomrule
\end{tabular}
\caption*{\textit{\textbf{Fuente.} Elaboración propia a partir de \path{INFO/CONVERSACIONES_IA}.}}
\end{table}

\section{Prompts depurados incorporados}

\newcommand{\promptDepurado}{\Needspace{7\baselineskip}\textbf{Prompt depurado.}}

{\small
\begin{enumerate}[leftmargin=*,itemsep=0.2\baselineskip,topsep=2pt,parsep=0pt,partopsep=0pt]
  \item \textbf{Prompt B.1. Análisis del estado del programa y mejora de parada}

  \textbf{Conversación base.} ChatGPT, mejora de parada precisa.

  \textbf{Objetivo.} Mejorar la precisión de parada ante señales de stop y semáforos sin depender de una recompensa genérica de supervivencia.

  \promptDepurado
  \begin{quote}
  \small
  Analiza las descripciones actuales del controlador de IA, señales de stop, semáforos, fitness y gestor evolutivo. Quiero que el modelo aprenda a detenerse cerca de \texttt{StopTargetLocation} cuando exista obligación de parar. Identifica qué información ya recibe la red, qué parte falta en el fitness y qué cambios son necesarios. Evita constantes arbitrarias: si propones umbrales, indica de qué variable del simulador deben derivarse. Separa la propuesta en pasos verificables y añade qué métricas habría que exportar para comprobar si la mejora funciona.
  \end{quote}

  \textbf{Respuesta aprovechada.} Se aprovechó la idea de distinguir entre tener información de parada como entrada y tener presión explícita en el fitness para aproximarse correctamente a la línea.

  \textbf{Comprobación realizada.} Se revisaron las descripciones fechadas de red neuronal, fitness e intersecciones, además de los CSV con tiempos de espera y completaciones de stop.

  \textbf{Uso final en el proyecto.} La memoria recoge la validación de stops, la ventana de parada, el aprendizaje de umbrales y la necesidad de diferenciar espera correcta de estancamiento.

  \item \textbf{Prompt B.2. Penalización de retroceso y estancamiento injustificado}

  \textbf{Conversación base.} Gemini, organización del proyecto, y ChatGPT, mejora de parada precisa.

  \textbf{Objetivo.} Evitar que el entrenamiento premie coches que sobreviven sin avanzar, retroceden o permanecen inmóviles sin una causa de tráfico válida.

  \promptDepurado
  \begin{quote}
  \small
  Revisa la lógica de fitness y detección de muerte para vehículos que retroceden o se quedan parados. Necesito diferenciar una parada legítima, por semáforo, stop, ceda u obstáculo delante, de una parada que solo alarga la vida del coche sin aportar recorrido útil. Propón una lógica basada en velocidad real, contexto de tráfico, sensores frontales y tiempo de vida. No uses números fijos sin justificación. Indica qué razón de muerte o métrica de diagnóstico debe registrarse para validar el cambio.
  \end{quote}

  \textbf{Respuesta aprovechada.} Se aprovechó la separación conceptual entre parada justificada y comportamiento perezoso, además de la necesidad de penalizar el retroceso con una causa trazable.

  \textbf{Comprobación realizada.} Se comprobó con los informes de entrenamiento, las razones de muerte \texttt{LAZY}, \texttt{BACKWARDS} y \texttt{TIMEFINISHED}, y los casos de coches vivos sin progreso.

  \textbf{Uso final en el proyecto.} El capítulo 6 describe la penalización por marcha atrás, estancamiento y navegación no válida, y el capítulo 7 advierte que el tiempo de vida no puede usarse como único indicador.

  \item \textbf{Prompt B.3. Diagnóstico de salida de calzada y coches bloqueados}

  \textbf{Conversación base.} Claude, análisis de problemas en datos de entrenamiento y comportamiento de vehículos.

  \textbf{Objetivo.} Explicar por qué algunos coches fuera de calzada o retenidos en señales llegaban a finalizar por tiempo en lugar de morir por una causa clara.

  \promptDepurado
  \begin{quote}
  \small
  Analiza el entrenamiento adjunto y las descripciones de fitness, intersecciones, gestor evolutivo y red neuronal. Busca coches que sobrevivan demasiado tiempo fuera de calzada o parados ante stops y cedas. Distingue hipótesis de evidencia en los CSV. Calcula qué condiciones deberían activar \texttt{NAV\_VIOLATION}, \texttt{LAZY}, \texttt{STOP\_VIOLATION} o \texttt{TIMEFINISHED}. Propón soluciones dinámicas, basadas en variables existentes, y explica cómo verificarlas con métricas exportadas.
  \end{quote}

  \textbf{Respuesta aprovechada.} Se aprovechó el análisis de interacción entre penalización acumulada, velocidad casi nula, contexto \texttt{ForceStop} y reinicio de contadores.

  \textbf{Comprobación realizada.} Se contrastó con registros de navegación fuera de calzada, finalizaciones por tiempo, penalizaciones acumuladas y ausencia de muertes \texttt{NAV\_VIOLATION} en determinadas sesiones.

  \textbf{Uso final en el proyecto.} La memoria incorpora la salida de calzada como problema de acumulación temporal y de penalización, y no como simple colisión geométrica.

  \item \textbf{Prompt B.4. Actualización del analizador de CSV}

  \textbf{Conversación base.} ChatGPT, ajuste del código de análisis y gráficas de tiempo vivo.

  \textbf{Objetivo.} Mantener los scripts de análisis alineados con las columnas nuevas exportadas por los Blueprints.

  \promptDepurado
  \begin{quote}
  \small
  Actualiza el script de análisis para los CSV actuales. Lee primero las cabeceras de \texttt{Fitness\_Debug.csv}, \texttt{Training\_Summary.csv} y \texttt{Test\_Summary.csv}. Añade compatibilidad con las nuevas columnas de stop, ceda, suavizado de dirección, genealogía, rotondas y solapes entre vehículos. El script debe validar cobertura, duplicados, filas incompletas, distribución de razones de muerte, evolución de fitness, vida media y diferencias por spawn. Si una columna no existe en una sesión antigua, debe degradar con aviso y no fallar.
  \end{quote}

  \textbf{Respuesta aprovechada.} Se aprovechó el enfoque de script robusto por cabeceras, con diagnóstico de cobertura y generación de informes comparables entre sesiones.

  \textbf{Comprobación realizada.} Se ejecutaron análisis sobre sesiones de abril, mayo y junio, comprobando informes, gráficas y coherencia con los CSV originales.

  \textbf{Uso final en el proyecto.} El capítulo 7 se apoya en métricas calculadas a partir de estos CSV y conserva explícitamente incidencias de cobertura o duplicados.

  \item \textbf{Prompt B.5. Separación de entrenamiento y test}

  \textbf{Conversación base.} ChatGPT, modificación de triggers y refresco de spawns.

  \textbf{Objetivo.} Evitar que el rendimiento de test se midiera sobre los mismos recorridos que condicionaban el entrenamiento.

  \promptDepurado
  \begin{quote}
  \small
  Revisa la lógica de spawns y triggers de dirección. Quiero separar entrenamiento y test: en entrenamiento deben usarse solo los puntos y trayectorias habilitados para aprender, mientras que el test debe evaluar un conjunto independiente y completo. Indica qué variables deben pertenecer al trigger, cuáles al gestor evolutivo y en qué momento conviene refrescar la lista de spawns para evitar mezclas entre modos. Propón comprobaciones en CSV para verificar que cada ejecución usa el conjunto correcto.
  \end{quote}

  \textbf{Respuesta aprovechada.} Se aprovechó la necesidad de separar \texttt{TrainMode}, listas válidas por modo, refresco de spawns y trayectorias permitidas.

  \textbf{Comprobación realizada.} Se revisaron \texttt{BP\_Trigger}, \texttt{RefreshSpawnPoints}, los resúmenes de test y la distribución de observaciones por los 54 spawns.

  \textbf{Uso final en el proyecto.} La memoria presenta la separación train/test como decisión metodológica clave y advierte que el test no debe interpretarse como nueva fase de aprendizaje.

  \item \textbf{Prompt B.6. Selección y deduplicación de modelos}

  \textbf{Conversación base.} ChatGPT, duplicados de modelos, y Claude, selección de modelo para reentrenamiento.

  \textbf{Objetivo.} Reducir pruebas repetidas sobre backups equivalentes y seleccionar candidatos de reentrenamiento con criterios trazables.

  \promptDepurado
  \begin{quote}
  \small
  Analiza los backups \texttt{.sav} de modelos y los tests asociados. Primero identifica modelos duplicados o casi idénticos comparando sus pesos, para no repetir evaluaciones. Después ordena los candidatos de reentrenamiento usando tasa de acierto, fitness medio, vida media, razones de muerte, estabilidad por spawn y coherencia entre entrenamiento y test. No elijas solo por \textit{BestFit}. Explica qué modelo conviene conservar como base, cuál usar como alternativa y qué pruebas faltan antes de aceptar el cambio.
  \end{quote}

  \textbf{Respuesta aprovechada.} Se aprovechó la idea de no tratar cada copia fechada como un modelo independiente si sus pesos no aportaban una política distinta.

  \textbf{Comprobación realizada.} Se compararon pesos, informes de test, métricas por sesión y diferencias entre \texttt{SuperCarModel} y \texttt{SessionCarModel}.

  \textbf{Uso final en el proyecto.} El capítulo 6 explica la persistencia de modelos y el capítulo 7 evita sustituir la referencia por sesiones posteriores no reproducibles.

  \item \textbf{Prompt B.7. Diagnóstico de rotondas e intersecciones}

  \textbf{Conversación base.} ChatGPT, revisión de comportamiento del coche y análisis de modelos.

  \textbf{Objetivo.} Analizar por qué el vehículo fallaba en rotondas, especialmente en salidas posteriores, y cómo ajustar la geometría o el fitness sin perder generalización.

  \promptDepurado
  \begin{quote}
  \small
  Revisa el test del modelo y las descripciones de intersecciones. Me interesa entender el comportamiento en rotondas: entrada, permanencia, elección de salida, aceleración, colisiones y completado. Separa problemas de geometría, escala del fitness y percepción por sensores. Si propones mover puntos de spline o modificar bordes de intersección, justifica cómo afectará a la trayectoria y qué métrica permitirá verificar la mejora. Incluye una revisión específica para primera, segunda y tercera o posteriores salidas.
  \end{quote}

  \textbf{Respuesta aprovechada.} Se aprovechó la distinción entre problema geométrico, problema de recompensa y problema de percepción contextual.

  \textbf{Comprobación realizada.} Se revisaron métricas de entrada y completado de rotonda, colisiones, saldo específico de rotonda y análisis por salida.

  \textbf{Uso final en el proyecto.} Los capítulos 6 y 7 describen las rotondas como subsistema propio, con penalización de aceleración, completaciones confirmadas y limitaciones de escala.

  \item \textbf{Prompt B.8. Análisis de regresión tras cambios de código}

  \textbf{Conversación base.} Claude, análisis de regresión y errores en stops, con atención a \texttt{STOP\_VIOLATION}.

  \textbf{Objetivo.} Entender por qué una modificación local podía empeorar el comportamiento global, especialmente en stops.

  \promptDepurado
  \begin{quote}
  \small
  Compara dos ejecuciones del mismo modelo antes y después de los cambios en intersecciones, fitness y red neuronal. Usa los informes y CSV para cuantificar la caída de rendimiento. Localiza qué razones de muerte aumentan, qué cambios de Blueprint coinciden temporalmente y qué hipótesis explican la regresión. No atribuyas causalidad sin prueba: separa correlación, evidencia directa y comprobaciones pendientes. Propón una secuencia de pruebas cortas en la que solo cambie un componente cada vez.
  \end{quote}

  \textbf{Respuesta aprovechada.} Se aprovechó la metodología de comparar el mismo modelo bajo condiciones próximas, localizar cambios de Blueprints y mirar razones de muerte antes de aceptar una mejora.

  \textbf{Comprobación realizada.} Se contrastaron los bloques de referencia y seguimiento de junio, con especial atención a \texttt{STOP\_VIOLATION}, \texttt{STOP\_VIOLATION\_TIMEOUT} y a la variabilidad entre sesiones.

  \textbf{Uso final en el proyecto.} La conclusión experimental indica que una mejora local del fitness puede degradar stops o generalización, por lo que la versión estable debe definirse por bloques reproducibles.

  \item \textbf{Prompt B.9. Mutación dinámica y persistencia de modelos}

  \textbf{Conversación base.} Claude, ratio dinámico basado en población, y ChatGPT, mejora del guardado de modelos.

  \textbf{Objetivo.} Mejorar la creación de generaciones y el guardado de modelos sin depender de proporciones fijas ni duplicar copias equivalentes.

  \promptDepurado
  \begin{quote}
  \small
  Revisa \texttt{StartGeneration}, \texttt{EndGeneration} y \texttt{SaveWeightsToSlot}. Quiero que la población se reparta de forma dinámica según \texttt{PopulationSize}, conservando un individuo sin mutar del mejor modelo histórico y otro del mejor modelo de sesión. El resto debe distribuirse entre mutaciones y candidatos aleatorios o derivados, sin números mágicos. Además, revisa el criterio de guardado para evitar duplicados de pesos, pero sin perder un modelo que tenga fitness o metadatos relevantes. Explica qué campos deben registrarse para auditar la decisión.
  \end{quote}

  \textbf{Respuesta aprovechada.} Se aprovechó el enfoque de ratios dependientes de población y la necesidad de guardar modelo, fitness y procedencia de forma coherente.

  \textbf{Comprobación realizada.} Se revisaron las descripciones de \texttt{EVOLUTIONMANAGER}, los backups de modelos y las métricas de genealogía por índice de coche.

  \textbf{Uso final en el proyecto.} La memoria describe las familias de mutación, el elitismo, \texttt{SessionCarModel}, \texttt{SuperCarModel} y la copia fechada del modelo anterior.

  \item \textbf{Prompt B.10. Organización del proyecto y reproducibilidad}

  \textbf{Conversación base.} Gemini, organización del proyecto, ChatGPT, limpieza de logs, y Gemini, errores de módulos en Unreal Engine.

  \textbf{Objetivo.} Mantener el proyecto reproducible pese a trabajar con Blueprints, binarios, logs extensos y resultados experimentales pesados.

  \promptDepurado
  \begin{quote}
  \small
  Organiza el repositorio y el flujo de trabajo de un proyecto con Blueprints, modelos guardados, logs, CSV de entrenamiento y memoria en LaTeX. Distingue archivos versionables, temporales y evidencia experimental; propón nombres de carpetas para reconstruir una sesión y un procedimiento de limpieza que conserve warnings, mensajes propios de Blueprints y datos útiles de depuración.
  \end{quote}

  \textbf{Respuesta aprovechada.} Se aprovechó la separación entre evidencia de proyecto, resultados generados, logs temporales y documentación final.

  \textbf{Comprobación realizada.} Se revisó la estructura \texttt{INFO}, los scripts de análisis, los ficheros temporales de LaTeX y la limpieza de imágenes generadas no necesarias.

  \textbf{Uso final en el proyecto.} La memoria conserva un anexo de fuentes, un anexo de prompts, reglas de limpieza del repositorio y una distinción entre datos primarios y salidas generadas.
\end{enumerate}
}
